\documentclass[slideColor]{prosper}
\usepackage{amsmath,amsfonts,amsthm}
%\usepackage{epsfig}
\include{Qcircuit}


\newcommand{\bbC}{{\mathbb{C}}}
\newcommand{\C}{{\mathbb{C}}}
\newcommand{\F}{{\mathbb{F}}}
\newcommand{\N}{{\mathbb{N}}}
\newcommand{\Q}{{\mathbb{Q}}}
\renewcommand{\R}{{\mathbb{R}}}
\newcommand{\Z}{{\mathbb{Z}}}

\newcommand{\cA}{{\mathcal{A}}}
\newcommand{\cD}{{\mathcal{D}}}
\newcommand{\cH}{{\mathcal{H}}}
\newcommand{\cN}{{\mathcal{N}}}
\newcommand{\cP}{{\mathcal{P}}}
\newcommand{\cQ}{{\mathcal{Q}}}
\newcommand{\cS}{{\mathcal{S}}}
\newcommand{\cU}{{\mathcal{U}}}
\newcommand{\cV}{{\mathcal{V}}}

\newcommand{\tr}{\mathop{\mathrm{tr}}}
\newcommand{\rank}{\mathop{\mathrm{rank}}}
\newcommand{\poly}{\mathop{\mathrm{poly}}}
\newcommand{\E}{\mathop{\mbox{$\mathbf{E}$}}}
\newcommand{\schur}{\mathop{\mathrm{Schur}}}
\newcommand{\schursmall}[2]{{\mathrm{S}}_{#1,#2}}
\newcommand{\planch}{\mathop{\mathrm{Planch}}}
\newcommand{\planchsmall}[1]{{\mathrm{P}}_{#1}}
\renewcommand{\d}{{\mathrm{d}}}

\newcommand{\HSP}{\textsc{hsp}\xspace}

\renewcommand{\>}{\rangle}
\newcommand{\<}{\langle}
%\newcommand{\ket}[1]{|#1\rangle}
%\newcommand{\bra}[1]{\langle #1|}
\newcommand{\scalar}[2]{\langle #1|#2\rangle}
\newcommand{\proj}[1]{\left|#1\right\>\!\left\<#1\right|}
\newcommand{\norm}[1]{\|#1\|}
\newcommand{\ot}{\otimes}
\newcommand{\eps}{\epsilon}
\newcommand{\ra}{\rightarrow}

\renewcommand{\setminus}{\mathbin{-}}
\newcommand{\partitionof}{\mathbin{\vdash}}


\newcommand{\I}{\mathrm{I}}
\renewcommand{\H}{\mathrm{H}}
\renewcommand{\S}{\mathrm{S}}
\newcommand{\T}{\mathrm{T}}
\newcommand{\G}{\mathrm{G}}
%

\newcommand{\bI}{\mathbb{I}}
\newcommand{\bH}{\mathbb{H}}
\newcommand{\bS}{\mathbb{S}}
\newcommand{\bT}{\mathbb{T}}




%%%%%%%%%%%%%%%%%%%%%%%%%%%%%%%%%%%%%%%%%%%%%%%%%%%%%%%%%%%%

\title{Classical Problems Characterizing\\ the Power of Quantum Computing}
\author{Pawel M. Wocjan\\[1ex]School of Electrical Engineering and Computer Science\\ University of Central Florida\\Orlando}
\email{wocjan@cs.ucf.edu}

\begin{document}

\maketitle

%%%%%%%%%%%%%%%%%%%%%%%%%%%%%%%%%%%%%%%%%%%%%%%%%%%%%%%%%%%%%%%%%%%%%%%%

\begin{slide}{Joint work with}
\begin{itemize}
\item Dominik Janzing (School of EECS, University of Central Florida) 
\item Shengyu Zhang (IQI, Caltech)
\end{itemize}

Talk is based on
\begin{itemize}
\item A single-shot measurement of the energy of product states in a translation invariant spin chain can replace any quantum computation, \texttt{quant-ph/0710.1615}

\medskip
\item A Simple PromiseBQP-complete Problem, {\em Theory of Computing}, Volume 3, Article 4, 2007

\medskip
\item A PromiseBQP-complete String Rewriting Problem, \texttt{quant-ph/0705.1180}
\end{itemize}
\end{slide}

%%%%%%%%%%%%%%%%%%%%%%%%%%%%%%%%%%%%%%%%%%%%%%%%%%%%%%%%%%%%%%%%%%%%%%%%%

\begin{slide}{Challenges in Quantum Computer Science}

\begin{itemize}
\item designing new efficient quantum algorithms
\item	understanding the computational power and limitations of quantum computers in comparision to those of classical computers 
\end{itemize}

\end{slide}

%%%%%%%%%%%%%%%%%%%%%%%%%%%%%%%%%%%%%%%%%%%%%%%%%%%%%%%%%%%%%%%%%%%%%%%%%

\begin{slide}{Complexity Class PromiseBQP}

PromiseBQP is the class of promise problems $(\Pi_{{\rm YES}},\Pi_{{\rm NO}})$ that can be solved by a {\em uniform} family of quantum circuits $\{U_n\}$ in the following sense:

\bigskip
Let ${\bf x}$ be a binary string of length $n$. Then the application of $U_n$ to $|{\bf x},{\bf 0}\>$ produces the state
\[
U_n |{\bf x},{\bf 0}\> = \alpha_{{\bf x},0}|0\> \otimes |\psi_{{\bf x},0}\> + \alpha_{{\bf x},1}|1\> \otimes |\psi_{{\bf x},1}\>
\]
such that 
\begin{itemize}
\item $|\alpha_{{\bf x},1}|^2\ge 2/3$ if $x\in \Pi_{{\rm YES}}$
\item $|\alpha_{{\bf x},1}|^2\le 1/3$ if $x\in \Pi_{{\rm NO}}$
\end{itemize}

\end{slide}

%%%%%%%%%%%%%%%%%%%%%%%%%%%%%%%%%%%%%%%%%%%%%%%%%%%%%%%%%%%%%%%%%%%%%%%%%

\begin{slide}{Measurements of Observables}

\begin{itemize}
\item
a measurement of the obervable $A$ with spectral decomposition 
\[
A=\sum_{\lambda} \lambda \, Q_\lambda
\]
is a quantum process that

\medskip
\begin{itemize}
\item outputs the eigenvalue $\lambda$  

\medskip
\item with probability ${\rm tr}(Q_\lambda \, \rho)$
\end{itemize}

\medskip
when applied to an arbitrary quantum state $\rho$
\end{itemize}

\end{slide}

%%%%%%%%%%%%%%%%%%%%%%%%%%%%%%%%%%%%%%%%%%%%%%%%%%%%%%%%%%%%%%%%%%%%%%%%%

\begin{slide}{Remarks: Measurements of Observables}

\begin{itemize}
\item note that

\begin{itemize}

\medskip
\item the post-measurement state is irrelevant

\medskip
\item a quantum process with the property that the expectation of the generated values is equal to ${\rm tr}(A\,\rho)$ may not suffice
\end{itemize}
\end{itemize}

\end{slide}

%%%%%%%%%%%%%%%%%%%%%%%%%%%%%%%%%%%%%%%%%%%%%%%%%%%%%%%%%%%%%%%%%%%%%%%%%

\begin{slide}{Spectral measure}

\begin{itemize}
\item
the spectral measure induced by the observable $A$ and the state $\rho$ is the probability distribution on $\Omega=\R$ where
\[
\Pr(\lambda) = {\rm tr}(Q_\lambda \rho)
\]
we set $Q_\lambda = {\bf 0}$ for all $\lambda\not\in {\rm Spec}(A)$

\item
a measurement of $A$ in the state $\rho$ allows us to sample according the spectral measure
\end{itemize}

\end{slide}

%%%%%%%%%%%%%%%%%%%%%%%%%%%%%%%%%%%%%%%%%%%%%%%%%%%%%%%%%%%%%%%%%%%%%%%%%

\begin{slide}{Accuracy of Measurements}

We say that an approximate measurement of $A$ has 
\begin{itemize}
\item maximal error $\delta$ and
\item reliability $1-\epsilon$
\end{itemize}
iff 
\[
\Pr(\lambda\in [a-\delta,b+\delta]) \ge {\rm tr}(Q_{[a,b]}\, \rho)\cdot(1-\epsilon)
\]
for all intervals $[a,b]\subseteq\R$ and all states $\rho$

\bigskip
where $Q_{[a,b]} = \sum_{\mu\in [a,b]} Q_\mu$

\end{slide}

%%%%%%%%%%%%%%%%%%%%%%%%%%%%%%%%%%%%%%%%%%%%%%%%%%%%%%%%%%%%%%%%%%%%%%%%%

\begin{slide}{Efficient Implementation of Measurements}
\begin{itemize}
\item let $A$ be a sparse matrix of size $N$

\medskip
\item a measurement of $A$ with maximal error $\delta$ and reliability $1-\epsilon$ can be implemented with resources (time and space) that are polynomial in $1/\delta$, $1/\epsilon$ and $\log N$

\medskip
\item the idea is to

\begin{itemize}
\item simulate time evolution $U\approx \exp(-i A/\|A\|)$ 

\medskip
\item apply phase estimation to $U$
\end{itemize}
\end{itemize}
\end{slide}

%%%%%%%%%%%%%%%%%%%%%%%%%%%%%%%%%%%%%%%%%%%%%%%%%%%%%%%%%%%%%%%%%%%%%%%%%

\begin{slide}{Quantum Computution with Measurements}
\begin{itemize}
\item there is a family of Hamiltonians $H^{(m)}$ acting on qudit chains $(\C^{60})^{\otimes m}$
\[
H^{(m)}=\sum_{j=0}^{m-2} H_{j,j+1}
\]
where $m=poly(n)$ such that a single-shot measurement of the energy of 

\begin{itemize}
\item canonical basis states 

\medskip
\item with error $\delta=1/poly(m)$ and reliability $1-\epsilon$ 
\end{itemize}

solves BQP in a probabilistic sense
\end{itemize}

\end{slide}

%%%%%%%%%%%%%%%%%%%%%%%%%%%%%%%%%%%%%%%%%%%%%%%%%%%%%%%%%%%%%%%%%%%%%%%%%

\begin{slide}{Properties of the Hamiltonians}

\begin{itemize}
\item the Hamiltonians $H^{(m)}$ are ``physically realistic'' since

\begin{itemize}
\item each $H^{(m)}$ is translation invariant

\medskip
\item each $H^{(m)}$ consists of nearest-neighbor interactions; $H$ acts on $(\C^{60})^{\otimes 2}$
\end{itemize}
\end{itemize}
\end{slide}

%%%%%%%%%%%%%%%%%%%%%%%%%%%%%%%%%%%%%%%%%%%%%%%%%%%%%%%%%%%%%%%%%%%%%%%%%

\begin{slide}{Solving BQP in a probabilistic sense}

\begin{itemize}
\item for all input strings ${\bf x}$ of length $n$ there is
  
  \begin{itemize}
  \item a partition $\R = {\cal Y}_m \, \dot{\cup} \, {\cal N}_m$ and

  \medskip
  \item a canonical basis state $|\psi_x\>\in(\C^{60})^{\otimes m}$ such that
  \end{itemize}
  \vspace{-0.5cm}
  \begin{eqnarray*}
  \Pr(\lambda\in \textcolor{red}{{\cal Y}_m}  \, | \, x\in \textcolor{red}{\Pi_{\rm YES}}) & \ge & p_{{\bf x}, \textcolor{red}{1}}  (1-\epsilon) \\ && \\
  \Pr(\lambda\in \textcolor{blue}{{\cal N}_m} \, | \, x\in \textcolor{blue}{\Pi_{\rm NO}}) & \ge & p_{{\bf x}, \textcolor{blue}{0}} (1-\epsilon) 
  \end{eqnarray*}

  where $\lambda$ is the random output of the approximate energy measurement of $H^{(m)}$ in the state $|\psi_x\>$
  
  \medskip
  \item
  the objects $m$, $|\psi_x\>$, $Y_m$, and $N_m$ can be efficiently computed from $x$ and $U^{(n)}$
\end{itemize}

\end{slide}

%%%%%%%%%%%%%%%%%%%%%%%%%%%%%%%%%%%%%%%%%%%%%%%%%%%%%%%%%%%%%%%%%%%%%%%%%

\begin{slide}{Idea of the Proof}
\begin{itemize}
\item let $U=U_r U_{r-1} \cdots U_1$ be an arbitrary quantum circuit 

\medskip
\item define the non-unitary gates $R:=|1\>\<1|$ and $E:={\bf 0}$ 

\medskip
\item for $s>r$ define the non-unitary circuit 
\[
\hat{U}=\hat{U}_{s+1}\hat{U}_s \hat{U}_{s-1} \cdots \hat{U}_{r+1} \hat{U}_r \hat{U}_{r-1} \cdots \hat{U}_1
\]
\begin{itemize}
\item the gates of $\hat{U}$ are 
\begin{eqnarray*}
\hat{U}_t     & := & U_t \quad\quad\mbox{for $t=1,\ldots,r$}    \\
\hat{U}_{r+1} & := & R                                          \\
\hat{U}_t     & := & I   \quad\quad\mbox{ for $t=r+2,\ldots,s$} \\
\hat{U}_{s+1} & := & E
\end{eqnarray*}
\end{itemize}
\end{itemize}
\end{slide}

%%%%%%%%%%%%%%%%%%%%%%%%%%%%%%%%%%%%%%%%%%%%%%%%%%%%%%%%%%%%%%%%%%%%%%%%%

\begin{slide}{Example}

\[
\Qcircuit @C=1em @R=.7em {
\lstick{|x\>} & \multigate{1}{U_1} & \qw                & \multigate{1}{U_3} & \gate{U_4} & \gate{R} & \gate{I} & \gate{I} & \gate{I} & \gate{E} & \qw \\
\lstick{|0\>} & \ghost{W}          & \multigate{1}{U_2} & \ghost{U_3}        & \qw        & \qw      & \qw      & \qw      & \qw      & \qw      & \qw \\
\lstick{|0\>} & \qw                & \ghost{U_2}        & \qw                & \qw        & \qw      & \qw      & \qw      & \qw      & \qw      & \qw  
} 
\]
\begin{itemize}
\item $r=4$ and $s=8$
\item read-out gate $R:=|1\>\<1|$ and end-gate $E:={\bf 0}$
\item 
\end{itemize}

\end{slide}

%%%%%%%%%%%%%%%%%%%%%%%%%%%%%%%%%%%%%%%%%%%%%%%%%%%%%%%%%%%%%%%%%%%%%%%%%

\begin{slide}{Forward-Time Operator}

\begin{itemize}
\item adjoin a Hilbert space $\cH_{\rm clock}$ to the quantum register 

\medskip
\item define the forward operator $F$ by
\begin{eqnarray*}
F \Big( |\varphi\> \otimes |t\> \Big)         & = & \hat{U}_{t+1} |\varphi\> \otimes |t+1\> \mbox{ for $t=0,\ldots,s$} \\
F^\dagger \Big( |\varphi\> \otimes |0\> \Big) & = & 0 
\end{eqnarray*}

\item define the initial state $|\psi\>:=|{\bf x},{\bf 0}\> \otimes |0\>_{\rm clock}$ 

\end{itemize}

\end{slide}

%%%%%%%%%%%%%%%%%%%%%%%%%%%%%%%%%%%%%%%%%%%%%%%%%%%%%%%%%%%%%%%%%%%%%%%%%

\begin{slide}{Intuitive Action of $F$}

\begin{itemize}
\item $F$ operates as follows:

\begin{itemize}
\item $F$ implements $U=U_r U_{r-1} \cdots U_1$

\medskip
\item $F$ checks the output of $U$ (with read-out gate $R$)

\medskip
\item \quad output NO $\Rightarrow$ $F$ stops

\medskip
\item \quad output YES $\Rightarrow$ $F$ carries out idle steps

\medskip
\item $F$ stops unconditionally (end-gate $E$)
\end{itemize}
\end{itemize}

\end{slide}

%%%%%%%%%%%%%%%%%%%%%%%%%%%%%%%%%%%%%%%%%%%%%%%%%%%%%%%%%%%%%%%%%%%%%%%%%

\begin{slide}{Action of $F$}

\begin{itemize}

\item define the matrix $\hat{F}$ acting on $\cH:=\C^{r+1}\oplus \C^{s+1}$ by 
\[
\hat{F}:=S_{r+1} \oplus S_{s+1}
\]
$S_{r+1}$ and $S_{s+1}$ are the adjacency matrices of line graphs with $r+1$ and $s+1$ vertices

\medskip
\item define the state $|\phi\>:=\sqrt{1-p_1} |0\> \oplus \sqrt{p_1} |0\>$

\medskip
\item $\Rightarrow$ there is a nice isomorphism between $F$ and $\hat{F}$
\begin{eqnarray*}
{\rm span}\{ F^t|\psi\>\} & \cong & \cH \\ 
&& \\
F^t|\psi\>                & \cong & \hat{F}^t|\phi\>
\end{eqnarray*}
\end{itemize}

\end{slide}

%%%%%%%%%%%%%%%%%%%%%%%%%%%%%%%%%%%%%%%%%%%%%%%%%%%%%%%%%%%%%%%%%%%%%%%%%

\begin{slide}{Spectral measure}

\begin{itemize}
\item $\Rightarrow$ the spectral measure induced by 
\[
H=F+F^\dagger \mbox{ and } |\psi\>=|{\bf x},{\bf 0}\>\otimes |0\>_{\rm clock}
\]
is a convex combination 
\[
(1-p_1) P_{r+1} + p_1 P_{s+1}
\]

where $P_{r+1}$ and $P_{s+1}$ are probability distributions on the spectra of the line graphs with $r+1$ and $s+1$ vertices

\medskip
\item the spectra ${\cal S}_{r+1}$ and ${\cal S}_{s+1}$ are disjoint and the distance between them is $\Omega(1/poly(r))$

\medskip 
\item 
$\Rightarrow$ quantum computation can be replaced by a single-shot approximate energy measurement
\end{itemize}
\end{slide}

%%%%%%%%%%%%%%%%%%%%%%%%%%%%%%%%%%%%%%%%%%%%%%%%%%%%%%%%%%%%%%%%%%%%%%%%%

\begin{slide}{A Simple PromiseBQP-complete Matrix Problem}

\begin{itemize}
\item let $A$ be a sparse matrix of size $N$

\medskip
\item given 
$m=polylog(N)$,
$j\in [1,\ldots,N]$,
$\epsilon=1/polylog(N)$, and 
$g\in [ - \|A\|^m, \|A\|^m ]$ 

decide if
\begin{itemize}
\item $(A^m)>$

\item 
\end{itemize}
\item
\end{itemize}

\end{slide}

%%%%%%%%%%%%%%%%%%%%%%%%%%%%%%%%%%%%%%%%%%%%%%%%%%%%%%%%%%%%%%%%%%%%%%%%%

\begin{slide}{A PromiseBQP-complete String Rewriting Problem}

\begin{itemize}
\item let $\cA$ be a finite alphabet

\medskip
\item let $\sim$ be a symmetric relation on substrings of length at most $k$ (where $k$ is a constant)

\medskip
\item let $A$ be the adjacency matrix of the undirected graph $G=(V,E)$, where
\begin{itemize}

\medskip
\item the set of vertices $V$ is $\cA^r$ (strings of length $r$ over $\cA$)

\medskip
\item the set of edges $E$ corresponds to possible conversions between strings, that is,
\[
(avb,awb)\in E \quad \Longleftrightarrow \quad v\sim w\
\]
\end{itemize}
\end{itemize}
\end{slide}

%%%%%%%%%%%%%%%%%%%%%%%%%%%%%%%%%%%%%%%%%%%%%%%%%%%%%%%%%%%%%%%%%%%%%%%%

\begin{slide}{String Rewriting Problem}
\begin{itemize}
\item let $s$, $t$, and $t'$ be three strings in $\cA^r$
\item let 
\[
\Delta_{s,t,t'}(n):=A_{s,t}^n - A_{s,t'}^n
\]
be the difference between the number of possibilities to obtain $t$ from $s$ and those to obtain $t'$ from $s$ after exactly $n$ replacements
\item given the promise
\begin{itemize}
\item $|\Delta_{s,t,t'}(n)| \le c^n$ for all $n\in\N$ ({\em growth condition})
\item $|\Delta_{s,t,t'}(m)|\ge \epsilon c^m$
\end{itemize}
determine the sign of $\Delta_{s,t,t'}(m)$ where

$c>0$ is a constant, $m\in\N$ with $m=poly(r)$, and $\epsilon\in [0,1]$ with $\epsilon=poly(1/r)$
\end{itemize}
\end{slide}

%%%%%%%%%%%%%%%%%%%%%%%%%%%%%%%%%%%%%%%%%%%%%%%%%%%%%%%%%%%%%%%%%%%%%%%%

\begin{slide}{Motivation from Bioinformatics}
\begin{itemize}
\item strings $s$, $t$, and $t'$ correspond to DNA sequences 
\item mutations change substrings (possible mutations are specified by the relation $\sim$)
\item String rewriting problem:

Are there more possibilities to mutate from $s$ to $t$ or from $s$ to $t'$ if exactly $m$ mutations occur?
\end{itemize}
\end{slide}

%%%%%%%%%%%%%%%%%%%%%%%%%%%%%%%%%%%%%%%%%%%%%%%%%%%%%%%%%%%%%%%%%%%%%%%%

\begin{slide}{PromiseBQP-complete}
Theorem: String Rewriting is PromiseBQP-complete

\bigskip
The string rewriting problem is the most difficult problem that can be efficiently solved on a quantum computer.

\bigskip
Quantum computers are better at estimating differences of combinatorial quantities than classical computers (provided that BPP is not equal to BQP).
\end{slide}

%%%%%%%%%%%%%%%%%%%%%%%%%%%%%%%%%%%%%%%%%%%%%%%%%%%%%%%%%%%%%%%%%%%%%%%%%

\begin{slide}{Conclusions}

\end{slide}

%%%%%%%%%%%%%%%%%%%%%%%%%%%%%%%%%%%%%%%%%%%%%%%%%%%%%%%%%%%%%%%%%%%%%%%%%

\begin{slide}{``Spectral'' Quantum Algorithms}

\begin{itemize}
\item Approximate evaluation of Jones, HOMFLYPT, and Tutte polynomials (PromiseBQP-complete problems)

\medskip
\item Evaluation of NAND-Trees

\end{itemize}

\end{slide}

%%%%%%%%%%%%%%%%%%%%%%%%%%%%%%%%%%%%%%%%%%%%%%%%%%%%%%%%%%%%%%%%%%%%%%%%%

\end{document}

%%%%%%%%%%%%%%%%%%%%%%%%%%%%%%%%%%%%%%%%%%%%%%%%%%%%%%%%%%%%%%%%%%%%%%%%%
%%%%%%%%%%%%%%%%%%%%%%%%%%%%%%%%%%%%%%%%%%%%%%%%%%%%%%%%%%%%%%%%%%%%%%%%%
%%%%%%%%%%%%%%%%%%%%%%%%%%%%%%%%%%%%%%%%%%%%%%%%%%%%%%%%%%%%%%%%%%%%%%%%%
%%%%%%%%%%%%%%%%%%%%%%%%%%%%%%%%%%%%%%%%%%%%%%%%%%%%%%%%%%%%%%%%%%%%%%%%%
%%%%%%%%%%%%%%%%%%%%%%%%%%%%%%%%%%%%%%%%%%%%%%%%%%%%%%%%%%%%%%%%%%%%%%%%%
%%%%%%%%%%%%%%%%%%%%%%%%%%%%%%%%%%%%%%%%%%%%%%%%%%%%%%%%%%%%%%%%%%%%%%%%%
%%%%%%%%%%%%%%%%%%%%%%%%%%%%%%%%%%%%%%%%%%%%%%%%%%%%%%%%%%%%%%%%%%%%%%%%%
%%%%%%%%%%%%%%%%%%%%%%%%%%%%%%%%%%%%%%%%%%%%%%%%%%%%%%%%%%%%%%%%%%%%%%%%%
%%%%%%%%%%%%%%%%%%%%%%%%%%%%%%%%%%%%%%%%%%%%%%%%%%%%%%%%%%%%%%%%%%%%%%%%%
%%%%%%%%%%%%%%%%%%%%%%%%%%%%%%%%%%%%%%%%%%%%%%%%%%%%%%%%%%%%%%%%%%%%%%%%%
%%%%%%%%%%%%%%%%%%%%%%%%%%%%%%%%%%%%%%%%%%%%%%%%%%%%%%%%%%%%%%%%%%%%%%%%%
%%%%%%%%%%%%%%%%%%%%%%%%%%%%%%%%%%%%%%%%%%%%%%%%%%%%%%%%%%%%%%%%%%%%%%%%%
%%%%%%%%%%%%%%%%%%%%%%%%%%%%%%%%%%%%%%%%%%%%%%%%%%%%%%%%%%%%%%%%%%%%%%%%%
%%%%%%%%%%%%%%%%%%%%%%%%%%%%%%%%%%%%%%%%%%%%%%%%%%%%%%%%%%%%%%%%%%%%%%%%%
%%%%%%%%%%%%%%%%%%%%%%%%%%%%%%%%%%%%%%%%%%%%%%%%%%%%%%%%%%%%%%%%%%%%%%%%%
%%%%%%%%%%%%%%%%%%%%%%%%%%%%%%%%%%%%%%%%%%%%%%%%%%%%%%%%%%%%%%%%%%%%%%%%%
%%%%%%%%%%%%%%%%%%%%%%%%%%%%%%%%%%%%%%%%%%%%%%%%%%%%%%%%%%%%%%%%%%%%%%%%%
%%%%%%%%%%%%%%%%%%%%%%%%%%%%%%%%%%%%%%%%%%%%%%%%%%%%%%%%%%%%%%%%%%%%%%%%%
%%%%%%%%%%%%%%%%%%%%%%%%%%%%%%%%%%%%%%%%%%%%%%%%%%%%%%%%%%%%%%%%%%%%%%%%%
%%%%%%%%%%%%%%%%%%%%%%%%%%%%%%%%%%%%%%%%%%%%%%%%%%%%%%%%%%%%%%%%%%%%%%%%%
%%%%%%%%%%%%%%%%%%%%%%%%%%%%%%%%%%%%%%%%%%%%%%%%%%%%%%%%%%%%%%%%%%%%%%%%%
%%%%%%%%%%%%%%%%%%%%%%%%%%%%%%%%%%%%%%%%%%%%%%%%%%%%%%%%%%%%%%%%%%%%%%%%%
%%%%%%%%%%%%%%%%%%%%%%%%%%%%%%%%%%%%%%%%%%%%%%%%%%%%%%%%%%%%%%%%%%%%%%%%%
%%%%%%%%%%%%%%%%%%%%%%%%%%%%%%%%%%%%%%%%%%%%%%%%%%%%%%%%%%%%%%%%%%%%%%%%%
%%%%%%%%%%%%%%%%%%%%%%%%%%%%%%%%%%%%%%%%%%%%%%%%%%%%%%%%%%%%%%%%%%%%%%%%%
%%%%%%%%%%%%%%%%%%%%%%%%%%%%%%%%%%%%%%%%%%%%%%%%%%%%%%%%%%%%%%%%%%%%%%%%%
%%%%%%%%%%%%%%%%%%%%%%%%%%%%%%%%%%%%%%%%%%%%%%%%%%%%%%%%%%%%%%%%%%%%%%%%%
%%%%%%%%%%%%%%%%%%%%%%%%%%%%%%%%%%%%%%%%%%%%%%%%%%%%%%%%%%%%%%%%%%%%%%%%%
%%%%%%%%%%%%%%%%%%%%%%%%%%%%%%%%%%%%%%%%%%%%%%%%%%%%%%%%%%%%%%%%%%%%%%%%%
%%%%%%%%%%%%%%%%%%%%%%%%%%%%%%%%%%%%%%%%%%%%%%%%%%%%%%%%%%%%%%%%%%%%%%%%%
%%%%%%%%%%%%%%%%%%%%%%%%%%%%%%%%%%%%%%%%%%%%%%%%%%%%%%%%%%%%%%%%%%%%%%%%%
%%%%%%%%%%%%%%%%%%%%%%%%%%%%%%%%%%%%%%%%%%%%%%%%%%%%%%%%%%%%%%%%%%%%%%%%%
%%%%%%%%%%%%%%%%%%%%%%%%%%%%%%%%%%%%%%%%%%%%%%%%%%%%%%%%%%%%%%%%%%%%%%%%%

\begin{slide}{Promise Problem}

A promise problem $\Pi$ is a pair $(\Pi_{\rm YES},\Pi_{\rm NO})$ of non-intersecting sets, i.e., 
$\Pi_{\rm YES},\Pi_{\rm NO}\subseteq\{0,1\}^*$ and $\Pi_{\rm YES}\cap\Pi_{\rm NO}=\emptyset$.

\bigskip
The set $\Pi_{\rm YES}\cup\Pi_{\rm NO}$ is called the {\em promise}.

\bigskip
Standard ``language recognition'' problems can be cast as the special case in which that promise is the set of all strings $\Pi_{\rm YES}\cup\Pi_{\rm NO}=\{0,1\}^*$ ({\em trivial promise})

\end{slide}

%%%%%%%%%%%%%%%%%%%%%%%%%%%%%%%%%%%%%%%%%%%%%%%%%%%%%%%%%%%%%%%%%%%%%%%%%

\begin{slide}{Complexity Class PromiseBQP}

PromiseBQP is the class of promise problems $(\Pi_{\rm YES},\Pi_{\rm NO})$ that can be solved by a {\em uniform} family of quantum circuits $\{U_n\}$ in the following sense:

\bigskip
Let ${\bf x}$ be a binary string of length $n$. Then the application of $U_n$ to $|{\bf x},{\bf 0}\>$ produces the state
\[
U_n |{\bf x},{\bf 0}\> = \alpha_{{\bf x},0}|0\> \otimes |\psi_{{\bf x},0}\> + \alpha_{{\bf x},1}|1\> \otimes |\psi_{{\bf x},1}\>
\]
such that 
\begin{itemize}
\item $|\alpha_{{\bf x},1}|^2\ge 2/3$ if $x\in\Pi_{\rm YES}$
\item $|\alpha_{{\bf x},1}|^2\le 1/3$ if $x\in\Pi_{\rm NO}$
\end{itemize}

note that nothing is required with respect to inputs that violate the promise
\end{slide}

%%%%%%%%%%%%%%%%%%%%%%%%%%%%%%%%%%%%%%%%%%%%%%%%%%%%%%%%%%%%%%%%%%%%%%%%%

\begin{slide}{Complete Problems}
\begin{itemize}
\item no complete problems are known for complexity classes such as BPP, BQP, and MA
\item complete problems are only known for their promise versions 
\end{itemize}
\end{slide}

%%%%%%%%%%%%%%%%%%%%%%%%%%%%%%%%%%%%%%%%%%%%%%%%%%%%%%%%%%%%%%%%%%%%%%%%

\begin{slide}{String Rewriting Problem}

\begin{itemize}
\item let $\cA$ be a finite alphabet

\medskip
\item let $\sim$ be a symmetric relation on substrings of length at most $k$ (where $k$ is a constant)

\medskip
\item let $A$ be the adjacency matrix of the undirected graph $G=(V,E)$, where
\begin{itemize}

\medskip
\item the set of vertices $V$ is $\cA^r$ (strings of length $r$ over $\cA$)

\medskip
\item the set of edges $E$ corresponds to possible conversions between strings, that is,
\[
(avb,awb)\in E \quad \Longleftrightarrow \quad v\sim w\
\]
\end{itemize}
\end{itemize}
\end{slide}

%%%%%%%%%%%%%%%%%%%%%%%%%%%%%%%%%%%%%%%%%%%%%%%%%%%%%%%%%%%%%%%%%%%%%%%%

\begin{slide}{String Rewriting Problem}
\begin{itemize}
\item let $s$, $t$, and $t'$ be three strings in $\cA^r$
\item let 
\[
\Delta_{s,t,t'}(n):=A_{s,t}^n - A_{s,t'}^n
\]
be the difference between the number of possibilities to obtain $t$ from $s$ and those to obtain $t'$ from $s$ after exactly $n$ replacements
\item given the promise
\begin{itemize}
\item $|\Delta_{s,t,t'}(n)| \le c^n$ for all $n\in\N$ ({\em growth condition})
\item $|\Delta_{s,t,t'}(m)|\ge \epsilon c^m$
\end{itemize}
determine the sign of $\Delta_{s,t,t'}(m)$ where

$c>0$ is a constant, $m\in\N$ with $m=poly(r)$, and $\epsilon\in [0,1]$ with $\epsilon=poly(1/r)$
\end{itemize}
\end{slide}

%%%%%%%%%%%%%%%%%%%%%%%%%%%%%%%%%%%%%%%%%%%%%%%%%%%%%%%%%%%%%%%%%%%%%%%%

\begin{slide}{Motivation from Bioinformatics}
\begin{itemize}
\item strings $s$, $t$, and $t'$ correspond to DNA sequences 
\item mutations change substrings (possible mutations are specified by the relation $\sim$)
\item String rewriting problem:

Are there more possibilities to mutate from $s$ to $t$ or from $s$ to $t'$ if exactly $m$ mutations occur?
\end{itemize}
\end{slide}

%%%%%%%%%%%%%%%%%%%%%%%%%%%%%%%%%%%%%%%%%%%%%%%%%%%%%%%%%%%%%%%%%%%%%%%%

\begin{slide}{PromiseBQP-complete}
Theorem: String Rewriting is PromiseBQP-complete

\bigskip
The string rewriting problem is the most difficult problem that can be efficiently solved on a quantum computer.

\bigskip
Quantum computers are better at estimating differences of combinatorial quantities than classical computers (provided that BPP is not equal to BQP).
\end{slide}

%%%%%%%%%%%%%%%%%%%%%%%%%%%%%%%%%%%%%%%%%%%%%%%%%%%%%%%%%%%%%%%%%%%%%%%%

\begin{slide}{String Rewriting is in PromiseBQP}

\begin{itemize}
\item rewrite $\Delta_{s,t,t'}(m)$ as a difference of expected values
\[
\Delta_{s,t,t'}(m)=\sqrt{2}(\<\psi^+|A^m|\psi^+\> - \<\psi^-|A^m|\psi^-\>)
\]
with $|\psi^{\pm}\>=(|s\>\pm|\varphi\>)/\sqrt{2}$ and $|\varphi\>=(|t\>+|t'\>))/\sqrt{2}$
\item determine the expected values as follows:
\begin{itemize}
\item let $U:=\exp(-i A/b)$ where $b\le\|A\|$
\item apply quantum phase estimation to $U$ and $|\psi^{\pm}\>$

\medskip
this allows one to sample according to the spectral measure induced by $A$ and $|\psi^{\pm}\>$: 

eigenvalue $\lambda_j$ occurs with probability  $\<\psi^{\pm}|P_j|\psi^{\pm}\>$
\end{itemize}
\item raise the outcome to the $m$th power
\item output average value after sufficiently many runs 
\end{itemize}
\end{slide}

%%%%%%%%%%%%%%%%%%%%%%%%%%%%%%%%%%%%%%%%%%%%%%%%%%%%%%%%%%%%%%%%%%%%%%%%

\begin{slide}{Technical Difficulties}

\begin{itemize}
\item imperfect simulated time evolution
\item imperfect reproduction of the spectral measure
\item amplification of errors due to $x\mapsto x^m$ 

$\Longrightarrow$ use a Lipschitz-continuous function instead

\item imperfect estimation due to statistical fluctuations 

\item most serious problem: the accuracy of the estimation is of order $b^m$  

$\Longrightarrow$ the growth condition allows one to zoom onto the right scale, making it possible to attain accuracy of order $c^m$ 
\end{itemize}
\end{slide}

%%%%%%%%%%%%%%%%%%%%%%%%%%%%%%%%%%%%%%%%%%%%%%%%%%%%%%%%%%%%%%%%%%%%%%%%

\begin{slide}{String Rewriting is PromiseBQP-hard}
\end{slide}

%%%%%%%%%%%%%%%%%%%%%%%%%%%%%%%%%%%%%%%%%%%%%%%%%%%%%%%%%%%%%%%%%%%%%%%%

\begin{slide}{Equivalent Definition of PromiseBQP}

\begin{itemize}
\item let $U$ be a quantum circuit and let ${\bf x}$ be the input
\item let $p_0$ and $p_1$ denote the probability of finding the output qubit in state $|0\>$ and $|1\>$, respectively
\item elementary calculations show that
\[
\<{\bf x}\,, 0\ldots 0|U^\dagger\sigma_z U|{\bf x}\,, 0\ldots 0\>=p_0 - p_1
\]
\item the sign of this quantity indicates if $x\in\Pi_{\rm YES}$ or $x\in\Pi_{\rm NO}$
\end{itemize}

\end{slide}

%%%%%%%%%%%%%%%%%%%%%%%%%%%%%%%%%%%%%%%%%%%%%%%%%%%%%%%%%%%%%%%%%%%%%%%%

\begin{slide}{Quantum-Universal Interaction}
we construct a quantum-universal translationally-invariant finite-range interaction on a line 
\[
H = F + F^\dagger = \sum_j V_j + \sum_j V_j^\dagger
\]
acting on $\cH:=(\C^d)^{\otimes l}$ such that
\[
\<\omega|H^m|\alpha\>=(p_0 - p_1) \cdot w(m)
\]

where $|\alpha\>$ and $|\omega\>$ are basis states of $\cH$ that encode the input $x$ and the circuit $U$
and $w$ is a certain function with $w : \N \rightarrow \R_+$

\bigskip 
$\Longrightarrow$ the sign of $\<\omega|H^m|\alpha\>$ indicates if $x\in\Pi_{\rm YES}$ or $x\in\Pi_{\rm NO}$
\end{slide}

%%%%%%%%%%%%%%%%%%%%%%%%%%%%%%%%%%%%%%%%%%%%%%%%%%%%%%%%%%%%%%%%%%%%%%%%

\begin{slide}{Action of Forward-Time Operator}
\[
\begin{array}{cccccccccccccccccc}
\#   &\#      &\#      &\#      &\#      &\boxplus & \I  & \T  & \I  &\H      &\I      &\H      &\S      &\I  &\#     &\#  \\
\mid &\bullet &\bullet &\bullet &\bullet &\mid     & b_1 & b_2 & b_3 & b_4    &\bullet &\bullet &\bullet &\bullet &\bullet &\bullet
\end{array}\,
\]
\end{slide}

%%%%%%%%%%%%%%%%%%%%%%%%%%%%%%%%%%%%%%%%%%%%%%%%%%%%%%%%%%%%%%%%%%%%%%%%

\begin{slide}{Action of Forward-Time Operator}
\[
\begin{array}{cccccccccccccccccc}
\#   &\#      &\#      &\#      &\#      &\blacksquare & \I  & \T  & \I  &\H      &\I      &\H      &\S      &\I  &\#     &\#  \\
\mid &\bullet &\bullet &\bullet &\bullet &\mid    & b_1 & b_2 & b_3 & b_4    &\bullet &\bullet &\bullet &\bullet &\bullet &\bullet
\end{array}\,
\]
\end{slide}

%%%%%%%%%%%%%%%%%%%%%%%%%%%%%%%%%%%%%%%%%%%%%%%%%%%%%%%%%%%%%%%%%%%%%%%%

\begin{slide}{Action of Forward-Time Operator}
\[
\begin{array}{cccccccccccccccccc}
\#   &\#      &\#      &\#      &\#      &\I      & \blacksquare  & \T  & \I  &\H      &\I      &\H      &\S      &\I  &\#     &\#  \\
\mid &\bullet &\bullet &\bullet &\bullet &\mid    & b_1 & b_2 & b_3 & b_4    &\bullet &\bullet &\bullet &\bullet &\bullet &\bullet
\end{array}\,
\]
\end{slide}

%%%%%%%%%%%%%%%%%%%%%%%%%%%%%%%%%%%%%%%%%%%%%%%%%%%%%%%%%%%%%%%%%%%%%%%%

\begin{slide}{Action of Forward-Time Operator}
\[
\begin{array}{cccccccccccccccccc}
\#   &\#      &\#      &\#      &\#      &\I      & \T  & \blacksquare  & \I  &\H      &\I      &\H      &\S      &\I  &\#     &\#  \\
\mid &\bullet &\bullet &\bullet &\bullet &\mid    & b_1 & b_2 & b_3 & b_4    &\bullet &\bullet &\bullet &\bullet &\bullet &\bullet
\end{array}\,
\]
$\blacksquare$ moves to the end of the program code
\end{slide}

%%%%%%%%%%%%%%%%%%%%%%%%%%%%%%%%%%%%%%%%%%%%%%%%%%%%%%%%%%%%%%%%%%%%%%%%

\begin{slide}{Action of Forward-Time Operator}
\[
\begin{array}{cccccccccccccccccc}
\#   &\#      &\#      &\#      &\#      &\I      & \T  & \I  &\H      &\I      &\H      &\S      &\I   &\blacksquare &\#     &\#  \\
\mid &\bullet &\bullet &\bullet &\bullet &\mid    & b_1 & b_2 & b_3 & b_4    &\bullet &\bullet &\bullet &\bullet &\bullet &\bullet
\end{array}\,
\]
\end{slide}

%%%%%%%%%%%%%%%%%%%%%%%%%%%%%%%%%%%%%%%%%%%%%%%%%%%%%%%%%%%%%%%%%%%%%%%%

\begin{slide}{Action of Forward-Time Operator}
\[
\begin{array}{cccccccccccccccccc}
\#   &\#      &\#      &\#      &\#      &\I      & \T  & \I  &\H      &\I      &\H      &\S      &\I   &\Diamond &\#     &\#  \\
\mid &\bullet &\bullet &\bullet &\bullet &\mid    & b_1 & b_2 & b_3 & b_4    &\bullet &\bullet &\bullet &\bullet &\bullet &\bullet
\end{array}\,
\]
\end{slide}

%%%%%%%%%%%%%%%%%%%%%%%%%%%%%%%%%%%%%%%%%%%%%%%%%%%%%%%%%%%%%%%%%%%%%%%%

\begin{slide}{Action of Forward-Time Operator}
\[
\begin{array}{cccccccccccccccccc}
\#   &\#      &\#      &\#      &\#      &\I      & \T  & \I  &\H      &\I      &\H      &\S      &\bI   &\# &\#     &\#  \\
\mid &\bullet &\bullet &\bullet &\bullet &\mid    & b_1 & b_2 & b_3 & b_4    &\bullet &\bullet &\bullet &\bullet &\bullet &\bullet
\end{array}\,
\]
\end{slide}

%%%%%%%%%%%%%%%%%%%%%%%%%%%%%%%%%%%%%%%%%%%%%%%%%%%%%%%%%%%%%%%%%%%%%%%%

\begin{slide}{Action of Forward-Time Operator}
\[
\begin{array}{cccccccccccccccccc}
\#   &\#      &\#      &\#      &\#      &\I      & \T  & \I  &\H      &\I      &\H      &\bS      &\I   &\# &\#     &\#  \\
\mid &\bullet &\bullet &\bullet &\bullet &\mid    & b_1 & b_2 & b_3 & b_4    &\bullet &\bullet &\bullet &\bullet &\bullet &\bullet
\end{array}\,
\]
\end{slide}

%%%%%%%%%%%%%%%%%%%%%%%%%%%%%%%%%%%%%%%%%%%%%%%%%%%%%%%%%%%%%%%%%%%%%%%%

\begin{slide}{Action of $H$ on the Relevant Subspace}
\begin{itemize}
\item the restriction of the Hamiltonian $H$ to the span of $F^i|\alpha\>$ is unitarily equivalent to the adjaceny matrix of the line graph
\[
\cP:=\sum_{j=0}^{\ell-1} |j\>\<j+1| + |j+1\>\<j|
\]
where $\ell$ is the number of transitions until the program code has moved all the way to the left (so that no transition is possible anymore)
\item 
\[
\<\omega|H^m|\alpha\>=(p_0 - p_1) (\cP^m)_{0,\ell-1}
\]
\end{itemize}
\end{slide}

%%%%%%%%%%%%%%%%%%%%%%%%%%%%%%%%%%%%%%%%%%%%%%%%%%%%%%%%%%%%%%%%%%%

\begin{slide}{Hamiltonian $H$}
\[
H = F + F^\dagger = \sum_{j=1}^r V_j + V_j^\dagger
\]
the operator $V$ acts on three adjacent cells
\[
V = Z \big( U_\T P_\T + U_\S P_\S + U_\H P_\H + P_\I \big)
\]
\begin{itemize}
\item $Z$ defines the annihilation and creation of configurations 
\item $P_\G$ is a projector onto $|\G\>$ (indicating that gate $G$ has to be applied)
\item $U_\G$ implements the gate $\G$
\end{itemize}
\end{slide}

%%%%%%%%%%%%%%%%%%%%%%%%%%%%%%%%%%%%%%%%%%%%%%%%%%%%%%%%%%%%%%%%%%%

\begin{slide}{Conversion of $H$ into string rewriting rules}
\begin{itemize}
\item the entries of the Hamiltonian $H$ are $0,1,\pm 1/\sqrt{2}$
\item the entries should be only $0$ and $1$, so that the Hamiltonian defines an adjacency matrix
\item it is possible to construct a Hamiltonian $\tilde{H}$ such that it contains only entries $0,\pm 1/\sqrt{2}$
\item $\tilde{H}$ has the same spectral properties as $H$
\end{itemize}
\end{slide}

%%%%%%%%%%%%%%%%%%%%%%%%%%%%%%%%%%%%%%%%%%%%%%%%%%%%%%%%%%%%%%%%%%%

\begin{slide}{Conversion into Adjacency Matrix}
\begin{itemize}
\item the operator $\tilde{V}$ acts on three cells $(\C^d)^{\otimes 3}$
\[
\tilde{V}=\sum_{a,b} s_{ab}|a\>\<b|
\]
\item define operators $X_{ab}$ acting on $(\C^2)^{\otimes 3}$ (`minus-sign simulator'' cells) by setting
\begin{eqnarray*}
X_{ab} = I \otimes \sigma_x \otimes I &\mbox{ if }& s_{ab}=-1/\sqrt{2} \\
X_{ab} = I \otimes I        \otimes I &\mbox{ if }& s_{ab}=1/\sqrt{2}  \\
X_{ab} = {\bf 0}                      &\mbox{ if }& s_{ab}=0 
\end{eqnarray*}
\end{itemize}
\end{slide}

%%%%%%%%%%%%%%%%%%%%%%%%%%%%%%%%%%%%%%%%%%%%%%%%%%%%%%%%%%%%%%%%%%%%%

\begin{slide}{Conversion into Adjacency Matrix}
\begin{itemize}
\item prepare the simulator cells in the state
\[
|\psi\> = |0\> \otimes |-\> \otimes |0\> 
\]
\item then we can simulate the minus sign since
\[
X_{ab} |\psi\> = s_{ab} |\psi\>
\]
\item define the operator $\hat{V}$ by setting
\[
\hat{V}=\sqrt{2}\sum_{a,b} X_{ab} \otimes |a\>\<b|
\]
\item $\hat{V}$ defines the substitution rules for substrings
\item let $A$ be the corresponding Hamiltonian/adjacency matrix
\end{itemize}
\end{slide}

%%%%%%%%%%%%%%%%%%%%%%%%%%%%%%%%%%%%%%%%%%%%%%%%%%%%%%%%%%%%%%%%%%%%%%

\begin{slide}{Conversion into String Rewriting Problem}
\begin{itemize}
\item one can show that there are three strings $s:=\alpha_0$, $t:=\omega_0$, and $t':=\omega_1$ such that
\begin{eqnarray*}
(A^m)_{st} - (A^m)_{st'} & = & 
\<\omega_0|A^m|\alpha_0\> - \<\omega_1|A^m|\alpha_0\> \\
& = & 
(p_0 - p_1) \cdot \sqrt{2}^m \cdot (\cP^m)_{0,\ell-1}
\end{eqnarray*}

\item eigenvalues of the line graph are in $[-2,2]$

\medskip
$\Rightarrow$ growth condition is satisfied with $c:=2\sqrt{2}$

\item one can choose $m=poly(r)$ such that $\sqrt{2}^m \cdot (\cP^m)_{0,\ell-1} = 1/poly(r) \cdot c^m$

\medskip
$\Rightarrow$ it suffices to estimate $\Delta_{s,t,t'}(m)$ with precision $\epsilon c^m$ (where $\epsilon=1/poly(r)$) to determine the sign
\end{itemize}
\end{slide}

%%%%%%%%%%%%%%%%%%%%%%%%%%%%%%%%%%%%%%%%%%%%%%%%%%%%%%%%%%%%%%%%%%%%%%

\begin{slide}{Conclusions}
\begin{itemize}
\item a purely classical (quantum free) characterization of BQP
\item quantum computers are better at counting and comparing the number of solutions of combinatorial problems (provided that BQP is strictly larger BQP)
\end{itemize}
\end{slide}

\end{document}

